% IFN680 Project 6 technical report.
% Generated by tools/build_report.py from tools/report_template.tex. Every number below is
% printed by main_report.ipynb, except the design settings and the prior-work findings that
% tools/verify.py lists.
\documentclass[10pt,a4paper,twocolumn]{article}

\usepackage[margin=1.5cm,top=0.9cm,bottom=0.8cm]{geometry}
\usepackage{lmodern}
\usepackage[T1]{fontenc}
\usepackage[utf8]{inputenc}
\usepackage{microtype}
\usepackage{graphicx}
\usepackage{booktabs}
\usepackage{amsmath}
\usepackage[font=small,labelfont=bf,skip=3pt]{caption}
\usepackage{titlesec}
% Loaded last, as hyperref requires; it only fills the PDF's title and author fields.
\usepackage[hidelinks,pdftitle={Project 6: The Sharpness Quest},
            pdfauthor={Karan Rooprai, Nhu Hieu Nguyen}]{hyperref}

\titlespacing*{\section}{0pt}{4pt}{0.5pt}
\titleformat{\section}{\normalfont\bfseries}{\thesection.}{0.5em}{}
\setlength{\parskip}{0pt}
\setlength{\parindent}{1em}
\setlength{\columnsep}{0.6cm}
\setlength{\textfloatsep}{4pt}
\setlength{\dbltextfloatsep}{5pt}
\setlength{\intextsep}{3pt}
\setlength{\floatsep}{3pt plus 1pt minus 1pt}
\setlength{\tabcolsep}{4pt}
\pagestyle{empty}
% Leading 3% under the class default (11.6pt instead of 12pt on 10pt type): with it the five
% sections, both figures, Table 1 and the references share the two pages without cutting content.
\linespread{0.97}
% Never split a hyphenated word across a column or page break.
\brokenpenalty=10000

\renewcommand{\topfraction}{0.94}
\renewcommand{\dbltopfraction}{0.94}
\renewcommand{\bottomfraction}{0.6}
\renewcommand{\textfraction}{0.06}
\renewcommand{\floatpagefraction}{0.75}

\begin{document}

% The banner sits inside the \twocolumn optional argument, which is the only way a page-wide
% figure can appear on page 1 in the article class.
\twocolumn[{%
\begin{center}
{\large\bfseries Project 6: The Sharpness Quest. Does a Discriminator Loss Sharpen a cVAE?}\\[2pt]
{\bfseries Group 4}
\end{center}
\vspace{-2pt}
Student 1: Karan Rooprai \hfill Student 1 ID: n12498122 \\
Student 2: Nhu Hieu Nguyen \hfill Student 2 ID: n12194778 \\
\vspace{3pt}

\begin{center}
\includegraphics[width=0.84\textwidth]{figures/figure_1_grids.pdf}
\end{center}
\vspace{-8pt}
{\small\textbf{Figure 1:} Twelve samples per digit from the baseline cVAE (left) and the cVAE
with a discriminator loss (right). Both decode the \emph{same} twelve latent codes, one per column.
The classifier's accuracy on each grid is 1.0000 (left) and 0.9917 (right).\par}
\vspace{6pt}
}]

% A page-wide table can only float to the top of a later page, so it is declared here, on page 1,
% to appear at the top of page 2.
\begin{table*}[t]
\centering
\footnotesize
\begin{tabular}{@{}lrrrrll@{}}
\toprule
Measure (10,000 test images) & Real & Baseline & Disc.\ loss & $\Delta$ & 95\% CI & $\Delta$ at seeds 1, 2 \\
\midrule
Reconstruction MSE & n/a & 0.01109 & 0.01194 & $+$0.00085 & [$+$0.00080, $+$0.00090] & $+$0.00238, $+$0.00067 \\
Laplacian variance, reconstructions & 0.2277 & 0.1184 & 0.1931 & $+$0.0747 & [$+$0.0741, $+$0.0753] & $+$0.0728, $+$0.0810 \\
Mid-grey fraction, reconstructions & 0.1439 & 0.2084 & 0.1533 & $-$0.0551 & [$-$0.0556, $-$0.0547] & $-$0.0555, $-$0.0590 \\
Laplacian variance, samples & 0.2277 & 0.1141 & 0.2055 & $+$0.0913 & [$+$0.0904, $+$0.0923] & $+$0.0906, $+$0.0943 \\
Mid-grey fraction, samples & 0.1439 & 0.2113 & 0.1469 & $-$0.0644 & [$-$0.0651, $-$0.0637] & $-$0.0662, $-$0.0672 \\
\midrule
Classifier accuracy, samples & 0.9836 & 0.9855 & 0.9875 & \multicolumn{3}{l}{seeds: 0.9855 0.9861 0.9879 / 0.9875 0.9877 0.9911} \\
Fréchet distance, samples & 0.45 & 51.32 & 4.81 & \multicolumn{3}{l}{seeds: 51.32 56.71 52.93 / 4.81 7.14 4.43} \\
Fréchet distance, reconstructions & 0.45 & 38.62 & 5.71 & \multicolumn{3}{l}{seeds: 38.62 40.55 39.66 / 5.71 7.18 4.89} \\
Within-class spread, samples & 1 & 0.640 & 0.883 & \multicolumn{3}{l}{seeds: 0.640 0.639 0.633 / 0.883 0.865 0.879} \\
\bottomrule
\end{tabular}
\caption{Seed-0 models; $\Delta$ is discriminator loss minus baseline, with a paired bootstrap
interval. In the lower block ``Real'' is the real digits' accuracy, the real-against-real floor
and 1 by definition, and seeds read baseline / discriminator loss.}
\end{table*}

\section{Introduction}
A conditional variational autoencoder (cVAE; Kingma and Welling, 2014; Sohn et al., 2015) encodes
an image with its class into a latent code $z$ and decodes it back; decoding a draw from
$\mathcal{N}(0, I)$ under a chosen class generates a new digit. Trained on pixel-wise mean
squared error (MSE), it blurs: when several images fit one code, the MSE-optimal output is their
mean, a grey smear of differently placed strokes. A discriminator (Goodfellow et al., 2014), a
classifier trained to tell real images from generated ones, penalises any output that does not
look like data, such as that smear. Hence the hypothesis tested here: \emph{``adding a
discriminator term to the loss function of a cVAE during training increases the sharpness of the
reconstructed images''}.

Prior work supports it, mostly by eye. Larsen et al.\ (2016) trained a VAE jointly with a GAN
(VAE/GAN) and found its samples, and its reconstructions of held-out CelebA faces, sharper than a
plain VAE's. The conditional CVAE-GAN of Bao et al.\ (2017) produced images that ``look realistic
and clear'' where a CVAE's were blurry. Rosca et al.\ (2017) added adversarial terms to an
autoencoder ($\alpha$-GAN) and obtained better reconstructions than the AGE autoencoder, but
mixed sample quality. Dosovitskiy and Brox (2016) found that image-space losses
over-smooth, and an added adversarial term sharpens but can invent content. Most of those change
several things at once; here the discriminator term is the only change, and sharpness is measured.

\section{Methodology}
\textbf{Baseline cVAE.} Tutorial 8.3's model has a convolutional encoder and a mirrored decoder,
both conditioned on a one-hot class, and the brief's 7 latent dimensions; it trains on
$100\,\mathrm{MSE} + 0.1\,\mathrm{KL}$ with AdamW (learning rate 0.001) in batches of 1,000. For
the $20\times20$ images the encoder's last kernel and the decoder's first become $3\times3$, the
only architectural change (350,107 parameters). At the tutorial's constant rate the loss kept
falling long after its 50 epochs, so both models follow a cosine schedule from 0.001 to 0 over
150 epochs.

\textbf{Discriminator loss.} A conditional discriminator $D$ (two strided convolutions and a
linear layer, 23,297 parameters) sees the image with its class as ten constant channels. As
Salvado (2026, p.~41) prescribes, real images are labelled 1 and fakes 0, $D$ minimises binary
cross-entropy, and each batch takes a $D$ step, then a cVAE step. The fakes are reconstructions
$\mathrm{Dec}(z, c)$, which the hypothesis names, and prior samples $\mathrm{Dec}(z_p, c)$, which
Task 3 evaluates, at half weight each, as in Larsen et al. The cVAE minimises
$100\,\mathrm{MSE} + 0.1\,\mathrm{KL} + \lambda L_G$ with $L_G = -\log D(\mathrm{fake})$, the
non-saturating form of the brief's $\log(1 - D(G(z)))$: it shares the fixed point but keeps its
gradient while $D$ confidently rejects early fakes (Goodfellow et al., 2014, \S3). $L_G$ sees $z$
detached, so only the decoder learns to fool $D$, again as in Larsen et al. $D$ uses AdamW
(learning rate $2\times10^{-4}$) on the same schedule.

\textbf{Protocol.} A seeded random 10,000 of the 60,000 training images form a development split.
There, $\lambda$ was chosen from $\{0.03, 0.1, 0.3, 1.0\}$ as the value whose reconstructions came closest
to the real development images' Laplacian variance (0.2311), excluding runs where $D$ passed 98\%
accuracy. In that order, the four runs' reconstructions scored 0.1466, 0.1732, 0.1909 and 0.1857, with $D$ accuracies of 0.8749, 0.7299, 0.6331 and 0.5997, so $\lambda = 0.3$ was chosen. Training counts as converged when the mean development reconstruction
loss over the last 10 epochs is within 1\% of the previous 10; both runs meet it (relative changes of $-$0.0013 and $-$0.0012; Figure~2b).
The final models were retrained on all 60,000 images with seeds 0,~1 and~2; only
\texttt{main\_report.ipynb} scores the test split.

\section{Experiments}
Each of the 10,000 test images is reconstructed from its encoder mean, and one prior sample
per image is drawn under its label from the same code for both models, so all comparisons are
paired. \textbf{Sharpness} is the variance of the Laplacian (Pech-Pacheco et al., 2000), read
against the real images' own 0.2277, since noise raises it too. The \textbf{mid-grey
fraction} (pixels strictly between 0.1 and 0.9) is the check: a share of 0.700 of real
pixels is exactly white, so blur shows as grey, and noise would raise both measures together.
\textbf{Class consistency} uses a separate classifier trained on the training images
(0.9836 test accuracy), as Salvado (2026, pp.~42, 87) proposes. \textbf{Quality and
variety} use the Fr\'echet distance (Heusel et al., 2017) in that classifier's 128-dimensional
features, since Inception-v3 was built for $299\times299$ colour photographs. It compares
class-stratified test halves (4,999 and 5,001 images), so the real-against-real distance
is a fair floor, and within-class feature spread, generated over real, tests for mode collapse.
Differences are discriminator loss minus baseline, with 95\% paired bootstrap intervals (2,000
resamples) at seed 0, repeated at seeds 1 and 2. \textbf{The rule was fixed in advance}: an output
is sharper if its Laplacian variance rises with an interval above zero and at every seed, and its
mid-grey fraction falls.

\section{Discussion}
In Figure~1 every baseline digit is recognisable but soft: its strokes are mostly grey and fray into
halos, most visibly on the loops of 5, 6 and 8. With the discriminator loss the same codes decode to
crisp, clean-edged strokes on a white background, so sharpness and overall quality improve together. Class
consistency holds in both grids, and each column keeps its slant and style across all ten rows.
Variety is kept too, and stroke weight ranges more widely, from hairline to bold, as in real
handwriting; the within-class spread in Table~1 agrees.

For reconstructions the rule is met: Laplacian variance rises from 0.1184 to 0.1931 (Table~1; Figure~2a), with an interval above zero and differences of $+$0.0728 and $+$0.0810 at seeds 1 and 2, and the mid-grey fraction falls from 0.2084 to 0.1533. Both values remain below the real test images' 0.2277, so the change moves reconstructions toward real sharpness, not past it. The samples meet it too: Laplacian variance rises from 0.1141 to 0.2055 ($+$0.0906 and $+$0.0943 at seeds 1 and 2), and the mid-grey fraction falls from 0.2113 to 0.1469. The price is pixel fidelity: reconstruction MSE rises from 0.01109 to 0.01194, beyond evaluation noise. Pixel error punishes small shifts (Larsen et al.): a crisp stroke slightly misplaced costs more than a blur over both places.

Both models' samples are read as their intended class about as often as real digits are (0.9855 and 0.9875 against 0.9836). In the classifier's feature space the samples' Fréchet distance to real digits falls from 51.32 to 4.81, at every seed (the real-against-real floor is 0.45), and the reconstructions' from 38.62 to 5.71. Within-class spread is 0.640 for the baseline and 0.883 with the discriminator loss (1 matches real variety), so samples with the discriminator loss are more varied, not less.

\setcounter{figure}{1}
\begin{figure}[t]
\centering
\includegraphics[width=\columnwidth]{figures/figure_2_reconstructions.pdf}
\caption{(a) One test digit per class, rebuilt by each seed-0 model. (b) Development
reconstruction loss, with the discriminator's accuracy dashed.}
\end{figure}

\textbf{Strengths and limitations.} The comparison is controlled: the models share data,
architecture, optimiser, schedule, seeds and batch order, and the verdict rests on paired
intervals, three seeds and a rule stated in advance. But three seeds bound training variance
loosely (the MSE cost ranges from $+$0.00067 to $+$0.00238 across seeds), four values of $\lambda$ sample the trade-off coarsely,
Laplacian variance also rewards noise, one pixel of blur is a large part of a stroke at
$20\times20$, and the Fr\'echet distance uses one small classifier's features. The discriminator also costs time: an epoch took
5.40 s against 2.09 s.

\textbf{Recommendations.} Use the discriminator term where perceived sharpness matters more than
pixel accuracy, choosing $\lambda$ on held-out data against the real images' own sharpness rather
than maximising it, which rewards noise; keep the adversarial gradient out of the encoder and
anneal both optimisers, as here. Before generalising, add seeds, a finer $\lambda$ grid, a
perceptual measure from a second, larger network, and natural images.

\section{Conclusion}
The evidence supports the hypothesis: under a rule fixed before the results, adding a discriminator term made both the reconstructions and the samples of a cVAE measurably sharper, consistently across three training seeds. Pixel error rose slightly, class consistency was unchanged and variety grew. Next, Larsen et al.'s shift-tolerant, feature-wise reconstruction loss would test whether the remaining blur comes from pixel MSE itself.

\section*{References}
\footnotesize
\setlength{\parindent}{0pt}
\newcommand{\bib}[1]{\par\hangindent=1em #1}
\bib{Bao, J. et al. (2017). CVAE-GAN: fine-grained image
generation through asymmetric training. \emph{ICCV}.}
\bib{Dosovitskiy, A. and Brox, T. (2016). Generating images with perceptual similarity metrics
based on deep networks. \emph{NeurIPS}.}
\bib{Goodfellow, I. et al. (2014). Generative adversarial nets. \emph{NeurIPS}.}
\bib{Heusel, M. et al. (2017). GANs trained by a two time-scale update rule converge to a local
Nash equilibrium. \emph{NeurIPS}.}
\bib{Kingma, D. P. and Welling, M. (2014). Auto-encoding variational Bayes. \emph{ICLR}.}
\bib{Larsen, A. B. L. et al. (2016). Autoencoding
beyond pixels using a learned similarity metric. \emph{ICML}.}
\bib{Pech-Pacheco, J. L. et al. (2000). Diatom autofocusing in brightfield microscopy: a
comparative study. \emph{ICPR}.}
\bib{Rosca, M. et al. (2017). Variational approaches for auto-encoding generative adversarial
networks. arXiv:1706.04987.}
\bib{Salvado, O. (2026). Week 9: transforming distribution (GAN, diffusion, flow). IFN680
lecture, QUT.}
\bib{Sohn, K. et al. (2015). Learning structured output representation using deep
conditional generative models. \emph{NeurIPS}.}

\end{document}
